% Part: applied-modal-logics
% Chapter: epistemic-logic
% Section: bisimulations

\documentclass[../../../include/open-logic-section]{subfiles}

\begin{document}

\olfileid{aml}{el}{bsd}

\olsection{باي‌سيميولېشنونه}

د معرفتي ژبې د څرګندونې د ځواک په اړه يوه پاتې پوښتنه د مدلونو او
په هغو کښې د رښتينو فارمولونو تر منځ له اړيکې سره تړلې ده۔ د چوکاټ
د مطابقت له پايلو مو ليدلي چې کله په يوه چوکاټ کښې ځينې فارمولونه
معتبر وي، نو هغه چوکاټ ځينې خاصيتونه هم پوره کوي۔ خو آيا زموږ
مودالي ژبه، د بېلګې په توګه، هغه نړۍ چې انعکاسي غشے لري له
هغې د نړۍو له نامتناهي لړۍ څخه بېلولے شي چې هره نړۍ ئې بلې ته
لار لري؟ يعنې آيا کوم فارمول $A$ شته چې يوازې له دغو دوو نړۍو
څخه په يوه کښې رښتيا وي؟

باي‌سيميولېشن د اړيکيزو مدلونو تر منځ داسې اړيکه ده چې په وسيله
ئې وايو د هغو جوړښت په اغېزمنه معنا يو شان دے۔ لکه څنګه چې به
وګورو، دا د مدلونو د هم‌ارزۍ هغه معنا نيسي چې زموږ په معرفتي
ژبه کښې هم څرګندېدے شي۔

\begin{defn}[باي‌سيميولېشن]
فرض کړئ $M_1 = \tuple{W_1, R_1, V_1}$ او $M_2 = \tuple{W_2, R_2, V_2}$
دوه اړيکيز مدلونه وي۔ همدارنګه، $\mathcal{R} \subseteq W_1 \times W_2$
دې يوه دوه‌ځاييزه اړيکه وي۔ $\mathcal{R}$ هله
\emph{باي‌سيميولېشن} بولو چې د هر $\tuple{w_1, w_2} \in \mathcal{R}$
لپاره لاندې شرطونه وي:

\begin{enumerate}
  \item د ټولو !!{propositional variable}s~$p$ لپاره
  $w_1 \in V_1(p)$ هله او يوازې هله چې $w_2 \in V_2(p)$ وي۔
  \item د ټولو عاملانو $a \in G$ او نړۍو $v_1 \in W_1$ لپاره، که
  $R_{1_a} w_1 v_1$ وي، نو داسې $v_2 \in W_2$ شته چې
  $R_{2_a} w_2 v_2$ او $\tuple{v_1, v_2} \in
  \mathcal{R}$ وي۔
  \item د ټولو عاملانو $a \in G$ او نړۍو $v_2 \in W_2$ لپاره، که
   $R_{2_a} w_2 v_2$ وي، نو داسې $v_1 \in W_1$ شته چې
   $R_{1_a} w_1 v_1$ او $\tuple{v_1, v_2} \in
   \mathcal{R}$ وي۔
\end{enumerate}
\textbf{د ژباړې د نسخې څرګندونه (OLAML-003):} په منجمده سرچينه کښې
د دواړو شرطونو عاملان په A کښې شمېرل شوي، خو د دې څپرکي د
عاملانو سټ G تعريف شوے دے؛ دلته يوازې هماغه دوه نمايې له G
سره سمې شوې دي۔ نور فورمولونه او منجمده سرچينه عين دي۔

که د $M_1$ او $M_2$ تر منځ باي‌سيميولېشن نړۍوې $w_1$ او $w_2$
سره وتړي، نو $\tuple{M_1, w_1} \leftrightarroweq
\tuple{M_2, w_2}$ هم ليکلے شو، او $\tuple{M_1, w_1}$ او
$\tuple{M_2, w_2}$ \emph{باي‌سيميولر} بولو۔
\end{defn}

د باي‌سيميولېشن د اړيکې بېلابېلې مادې بېلابېل شيان تضمينوي۔
لومړۍ ماده ډاډ راکوي چې باي‌سيميولر نړۍوې به هماغه له موداليت
پرته فارمولونه پوره کوي، ځکه د ټولو !!{propositional variable}s
په اړه يې يو شان حالت ټاکي۔ نورې دوه مادې، چې په ترتيب سره کله
«مخ ته» او «شاته» بلل کېږي، ډاډ ورکوي چې د لاسرسي اړيکې به
يو شان جوړښت ولري۔

\begin{thm}
که $\tuple{M_1, w_1} \leftrightarroweq \tuple{M_2, w_2}$ وي، نو د
هر !!{formula}~$!A$ لپاره $\mSat{M_1}{!A}[w_1]$ هله او يوازې
هله چې $\mSat{M_2}{!A}[w_2]$ وي۔
\end{thm}

\begin{figure}
  \begin{center}
    \begin{tikzpicture}[modal]
      \node[world] (w1) {$w_1$}; 
      \node[world] (w2) [above left=of w1]{$w_2$}; 
      \node[world] (w3) [above right=of w1] {$w_3$};
      \draw[<->] (w1) to node {$a$} (w2);
      \draw[->,loop below] (w1) to node {$a$} (w1);
      \draw[<->] (w1) to [swap] node {$a$} (w3);
      \draw[->,loop above] (w2) to node {$a$} (w2) ;
      \draw[->,loop above] (w3) to node {$a$} (w3) ;
      
      \node[world](v1)[right of=w1, xshift=1.5in]{$v_1$};
      \node[world](v2)[above of=v1]{$v_2$};
      \draw[<->] (v1) to node {$a$} (v2);
      \draw[->,loop below] (v1) to node {$a$} (v1);
      \draw[->,loop above] (v2) to node {$a$} (v2) ;

      \draw[-,dotted] (w1) to (v1) ;
      \draw[-,dotted,bend left=45] (w2) to (v2) ;
      \draw[-,dotted,bend left=30] (w3) to (v2) ;
    \end{tikzpicture}
  \end{center}
  \caption{دوه باي‌سيميولر مدلونه۔}
  \ollabel{fig:bisimilar}
\end{figure}

که څۀ هم په \olref{fig:bisimilar} کښې ښودل شوي دوه مدلونه عين
يو شان نۀ دي، خو داسې باي‌سيميولېشن شته چې نړۍوې $w_1$ او~$v_1$
سره تړي۔ دا اړيکه $w_2$ او $w_3$ دواړه هم له~$v_2$ سره تړي؛
مطلب دا چې زموږ په مودالي ژبه کښې داسې څۀ نۀ شي څرګندېدے چې
په رښتيا يې له يو بل څخه بېل کړي۔ خو که $w_2$ او~$w_3$ بېلابېل
!!{propositional variable}s پوره کول، حالت به بدل و۔

\end{document}
